% The paper: the up-to-date version of the thesis, explained in words.
% Concepts only: formulas live in the handbook (cited by number with \hb),
% proofs in Lean (cited by name with \lean) and the blueprint.
% Outline, rules and status: paper/outline.md.
% Built by CI as a PDF (xelatex) and a searchable web page (plasTeX),
% published next to the handbook. scripts/paper_check.py checks every
% citation (run it with --fix after renumbering the handbook).
\documentclass[11pt]{article}

\usepackage{amssymb, amsmath}
\usepackage[margin=2.5cm]{geometry}
\usepackage{booktabs}
\usepackage[colorlinks=true, linkcolor=blue, urlcolor=blue]{hyperref}

\newcommand{\hmu}{\hat\mu}
\newcommand{\Var}{\operatorname{Var}}
\newcommand{\cov}{\operatorname{cov}}

% --- Citations (all checked by scripts/paper_check.py) ----------------------
\newcommand{\site}{https://mikaelandblom-coder.github.io/credible-lean}
\newcommand{\repo}{https://github.com/mikaelandblom-coder/credible-lean/blob/main}
% \lean{Name}: a proved Lean declaration, linked to its source line.
\newcommand{\lean}[1]{\href{\site/find/\#doc/#1}{\texttt{\detokenize{#1}}}}
% \hb{label}{n}: handbook Result n (label = its \label in handbook.tex).
\newcommand{\hb}[2]{\href{\site/handbook.pdf\#nameddest=result.#2}{Handbook~(#2)}}
% \hbsec{label}{n}: handbook Section n.
\newcommand{\hbsec}[2]{\href{\site/handbook.pdf\#nameddest=section.#2}{Handbook~\S#2}}
% \finding{n}: item n of FINDINGS.md.
\newcommand{\finding}[1]{\href{\repo/FINDINGS.md}{Finding~#1}}
% \file{path}: a file in the repository (a working note, a study script).
\newcommand{\file}[1]{\href{\repo/#1}{\texttt{\detokenize{#1}}}}
% \app{step}: a step or switch of the app.
\newcommand{\app}[1]{\href{\site/app/}{app: #1}}

% --- The section template (paper/outline.md) --------------------------------
\newcommand{\oneline}{\paragraph{In one sentence.}}
\newcommand{\idea}{\paragraph{The idea.}}
\newcommand{\watch}{\paragraph{When it applies, and what to watch out for.}}
\newcommand{\see}{\paragraph{See.}}
\newcommand{\notyet}[1]{\par\noindent\textit{Not yet written. Will cover: #1}\par}

% Working title (owner to confirm: paper/outline.md, "Open choices").
\title{Credibility from second moments: a guide}
\author{M.~Andblom}
\date{Draft}

\begin{document}
\maketitle

\begin{abstract}
  Credibility theory blends a segment's own experience with the portfolio's.
  This paper explains one framework for doing that which assumes nothing but
  means and covariances, lets segments be correlated, and contains classical
  Bühlmann--Straub as a special case. It is the up-to-date version of the
  master thesis it grew out of. It explains; it does not prove or list
  formulas. The formulas are in the
  \href{\site/handbook.pdf}{handbook}, cited by number, and every result is
  machine-checked in Lean, cited by name; the
  \href{\site/}{blueprint} has the proofs.
\end{abstract}

\tableofcontents

% Section 0 is "How to read this", so the numbers match paper/outline.md.
\setcounter{section}{-1}

% ===========================================================================
\section{How to read this}\label{sec:read}

\oneline This paper explains the ideas; the formulas, the proofs and the
working tool each live in a document of their own, and the paper points to
them instead of repeating them.

\idea The work comes in four documents, one for each kind of content.
\begin{itemize}
  \item \emph{The paper} (this one) explains. It is the up-to-date version of
    the master thesis it grew out of (Andblom, Stockholm University, 2023:2)
    and the first thing to read or cite. Each section is one concept: what
    it is, the idea in plain words, and what to watch out for. It has few
    formulas and no proofs.
  \item \emph{The \href{\site/handbook.pdf}{handbook}} holds the formulas:
    the key results, numbered, each with a short note on how to use it. The
    paper cites them as ``Handbook~(5)''.
  \item \emph{The \href{\site/}{blueprint}} is the proof ledger: every
    statement with its proof and a graph of what depends on what. You reach
    it by following a link, not by reading it front to back.
  \item \emph{The \href{\site/app/}{app}} shows the tools at work. You pick a
    world (cohorts, volumes, correlation, noise), draw data, and compare
    fitted models with the best blend for the true parameters. The fits see
    only the data.
\end{itemize}
Every section ends with a \emph{See} line: the handbook results, the Lean
names of the proofs, the app step to try, the findings it explains, and the
thesis sections it replaces. Appendix~\ref{app:changes} lists what changed
since the thesis, with a concordance from each thesis equation to its
handbook number and Lean name; Appendix~\ref{app:notation} has the notation
and its translation to Bühlmann--Straub.

\emph{What ``machine-checked'' means.} Every result the paper states is a
theorem in Lean, a proof assistant: a program that accepts a proof only if
each step follows from the previous ones and from the axioms of mathematics.
A Lean name such as \lean{Credibility.Core.bestZ_eq} links to the statement
and its proof. Continuous integration rebuilds all proofs on every change
and checks that none of them rests on an unfinished step or on an extra
axiom. So if the paper says a result is proved, the statement as written in
Lean is true; you do not need to trust the author, or the AI that wrote most
of the Lean code with the author, to believe it.

\emph{What it does not mean.}
\begin{itemize}
  \item \emph{The Lean statement is what is proved, not the paper's
    paraphrase.} The paper and the handbook state results in words and
    formulas; the hypotheses in full are in Lean. When the two seem to
    differ, Lean is right, and the paper has a mistake to fix.
  \item \emph{A theorem is about the model, not about your data.} That the
    best $z_j$ is optimal is proved; that your portfolio follows the model
    is an assumption you have to judge. The paper says where the model is
    fragile (each ``What to watch out for'').
  \item \emph{Numbers from simulations are not theorems.} Results from the
    app's studies hold in the worlds that were simulated; the paper says so
    and cites the script that produced them. Statements about the best
    blend with known parameters are exact formulas and need no simulation.
  \item \emph{Some arguments are on paper only.} Where a result is shown by
    hand but not in Lean, the paper says so (for example the converse in
    Section~\ref{sec:special}).
\end{itemize}

\emph{Findings.} Formalizing the thesis turned up gaps (a missing
hypothesis, a typo) and new results. They are numbered in a dated log,
\file{FINDINGS.md}, and cited as ``Finding~$n$''. The paper explains each
one where it belongs; Appendix~\ref{app:changes} collects them by thesis
section.

\watch
\begin{itemize}
  \item \emph{Equation numbers} such as ``eq.~(26)'' refer to the thesis as
    published; handbook numbers can change between versions, so cite the
    handbook with its version or, better, the Lean name, which is stable.
  \item \emph{Only means and covariances are assumed} throughout, unless a
    result says otherwise in its name (Section~\ref{sec:model}).
\end{itemize}

\see \hbsec{sec:setting}{1} ·
\lean{Credibility.Core.bestZ_eq} ·
\app{World} · \file{FINDINGS.md}, \file{paper/outline.md} · thesis §1.

% ===========================================================================
\section{What credibility is}\label{sec:credibility}

\oneline Credibility is mixing a cohort's own, noisy experience with the
steadier experience of the whole portfolio, in the proportion that makes the
mix as accurate as possible.

\idea Say you sell travel insurance through a dozen payment cards, and you
want the risk premium per cardholder for each card. A small card has three
years of data and a loss cost of 1.8 per cardholder; the portfolio as a
whole runs at 1.0. Which number do you charge?

Neither is right on its own. The card's own figure $\hmu_j$ aims at the right
target, but with so little data much of the 1.8 may be luck. The portfolio
figure $\hmu$ is stable, but it is blind to anything that really makes this
card different. Credibility charges something in between:
\[
  \hmu_j^{\,b} = z_j\,\hmu_j + (1 - z_j)\,\hmu .
\]
The \emph{credibility factor} $z_j$ is how much you trust the card's own
experience. At $z_j = 1$ you trust it fully; at $z_j = 0$ you charge the
portfolio rate. Everything in this paper is about choosing $z_j$, and the two
estimates it blends, well.

\emph{Well} means: the blend should miss the card's true risk premium $\mu_j$
by as little as possible on average. If both estimates are unbiased, missing
by little on average means a small error variance. That variance is a
parabola in $z_j$, so it has exactly one lowest point, the \emph{best}
credibility factor (\hb{res:bestZ}{1}). Where the vertex sits depends on
only six numbers: the variances of the two estimates, of the truth, and the
three covariances between them. Nothing else about the data matters, and the
result holds for any two unbiased estimates, not only the ones built in this
paper.

There is a second way to read the best $z_j$, and it is the one to keep in
mind. Look at the gap $\hmu_j - \hmu$ between the card and the portfolio.
Part of that gap is real (the card really is riskier) and part is noise. The
best $z_j$ is the regression coefficient of the truth on the gap: the share
of the gap you expect to be real (\hb{res:z}{5}). A big card with a
stable history has a gap that is mostly real, so $z_j$ is near one. A tiny
card's gap is mostly noise, so $z_j$ is near zero.

Building the blend takes three choices: how to average a card's years into
$\hmu_j$, how to average the cards into $\hmu$, and the factor $z_j$. They
can be made one at a time, because no choice changes what the others are
optimizing. Section~\ref{sec:weights} makes the first two,
Section~\ref{sec:z} the third.

This is the oldest idea in actuarial science that still has a name of its
own. Bühlmann (1967) gave it its modern form and Bühlmann and Straub (1970)
the version with volumes that most insurers use. What is new here is that
cohorts may be \emph{correlated}: a cold winter or a surge in travel moves
all the cards at once. Classical credibility assumes that away. This
framework does not, and it gets Bühlmann--Straub back when the correlation
is zero (Section~\ref{sec:special}).

\watch
\begin{itemize}
  \item \emph{Unbiased means: you cannot tell the cohorts apart in
    advance.} The model takes every cohort to have the same expected risk
    before you look at its data. If you know that one card covers more, or
    reports in another currency, correct for that first; otherwise the
    known difference looks like real spread and every $z_j$ comes out too
    high.
  \item \emph{The blend does not replace the cohort's own figure.} The
    historical $\hmu_j$ is still the right number for looking back. The
    blend is a forecast.
  \item \emph{With correlation, $z_j$ need not lie between 0 and 1.} When
    the portfolio shares a cohort's noise, trusting the cohort ``more than
    fully'' can be best. Section~\ref{sec:z} explains when, and what to do
    about it.
  \item \emph{The second moments are taken as given.} Here they are just
    numbers. Section~\ref{sec:moments} shows they really are the
    covariances of random variables, so the error variance is the true mean
    squared error, and Section~\ref{sec:estimation} how to estimate them.
\end{itemize}

\see \hb{res:bestZ}{1}, \hb{res:z}{5} ·
\lean{Credibility.BlendMoments.bestZ_optimal},
\lean{Credibility.BlendMoments.bestZ_unique},
\lean{Credibility.Core.bestZ_eq},
\lean{Credibility.Core.Estimates.bestZ_mse_optimal} ·
\app{One draw} · thesis §1, §2.1.

% ===========================================================================
\section{The model}\label{sec:model}

\oneline Each observation is the overall level, plus a cohort effect that
never changes, plus one period's noise; cohorts are linked only through the
noise of the same period, and only means and covariances are assumed.

\idea Write $\mu_{jt}$ for what cohort $j$ shows in period $t$: loss cost
per cardholder in underwriting year $t$, say. The model is
\[
  \mu_{jt} = \underbrace{\mu + \tau\,\eta_j}_{\mu_j} + \sigma_{jt}\,\varepsilon_{jt},
  \qquad
  \cov(\varepsilon_{it}, \varepsilon_{js}) = \rho_{ij}\,\delta_{ts}.
\]
Read it from left to right.
\begin{itemize}
  \item $\mu$ is the portfolio's overall level.
  \item $\eta_j$ is the \emph{cohort effect}: how far cohort $j$ really is
    from the rest, in units of $\tau$. It is drawn once and stays. The
    number $\tau^2$ says how different cohorts are; a homogeneous portfolio
    has small $\tau^2$. The cohort's true level is $\mu_j = \mu + \tau\eta_j$:
    the target.
  \item $\varepsilon_{jt}$ is the \emph{noise}: what one period of data gets
    wrong. Its size $\sigma_{jt}$ is small for a big cohort and large for a
    small one, typically $\sigma^2/w_{jt}$ for volume $w_{jt}$.
  \item $\rho_{ij}$ is the correlation of two cohorts' noise in the same
    period. A surge in travel raises claims on every card in that year, not
    equally. Noise in different periods is uncorrelated
    ($\delta_{ts}$ is 1 when $t = s$ and 0 otherwise).
\end{itemize}
Both $\eta_j$ and $\varepsilon_{jt}$ are centred with unit variance, and the
cohort effects are uncorrelated with each other and with the noise.

\emph{Why error terms?} Classical credibility gives each cohort a risk
parameter $\Theta_j$ and draws them independently. Independence is built
into that language, so saying that cohorts are correlated takes effort.
Written as error terms, the correlation is one line, the covariance above.
This is the random-effects notation statisticians use, and it reads the same
whether cohorts are correlated or not.

\emph{Only means and covariances.} The model says nothing about the shape of
the noise: no normal, no Poisson, no gamma. Everything in this paper, from
the best $z_j$ to the unbiased estimators, needs only first and second
moments, so it holds whatever the claims distribution. Results that do
assume a noise family will live apart, named as such
(Section~\ref{sec:open}).

\emph{Bühlmann--Straub inside.} In Bühlmann--Straub notation $\mu$ is the
collective premium $m$, $\tau^2$ is $a$, and $\sigma_{jt}^2$ is
$\sigma^2/w_{jt}$; the full translation is in \hbsec{sec:setting}{1}. Of the
classical independence assumptions, exactly one is dropped: that two
cohorts' observations in the same period are uncorrelated. Put
$\rho_{ij} = 0$ for $i \ne j$ and you are back in Bühlmann--Straub, with its
credibility factors exactly (\hb{res:bs}{9}).

\emph{Every cohort, every period.} Here each cohort is observed in the same
$T$ periods. Section~\ref{sec:unequal} lets cohorts enter and leave; this
model is its special case, and in Lean every result about it is proved
there and read off.

\watch
\begin{itemize}
  \item \emph{The correlations must fit together.} The table $(\rho_{ij})$
    has to be a genuine correlation matrix (positive semidefinite): three
    cohorts cannot each be strongly negatively correlated with the other
    two. If the table fails this, no random variables have these moments
    and every formula built on it is about nothing
    (\hb{res:exists}{21}; Section~\ref{sec:moments}).
  \item \emph{No correlation across periods.} A trend, a slow claims
    inflation, or a reporting pattern that carries one year's errors into
    the next breaks $\delta_{ts}$. That is an open direction
    (Section~\ref{sec:open}).
  \item \emph{$\tau^2$ is not what you see inside one cohort.} Looking at a
    single cohort's history, its true level is already fixed, so its spread
    over time shows only $\sigma_{jt}^2$. The spread between cohorts is what
    reveals $\tau^2$, which is why estimating it needs several cohorts
    (Section~\ref{sec:estimation}).
\end{itemize}

\see \hbsec{sec:setting}{1}, \hb{res:bs}{9}, \hb{res:exists}{21} ·
\lean{Credibility.bs_z},
\lean{Credibility.Model.exists_noise_iff} ·
\app{World} · thesis §2.3.

% ===========================================================================
\section{From data to \texorpdfstring{$K$}{K}}\label{sec:K}

\oneline Everything the framework gives you (the weights, the portfolio
estimate, every credibility factor) is read off one matrix $K$, the
covariance of the cohort estimates, provided $K$ is positive definite.

\idea Credibility needs second moments, and in this framework they all come
from one place. Collect the covariances of the cohort estimates in a matrix,
$K_{ij} = \cov(\hmu_i, \hmu_j)$. The recipe has three steps
(\hbsec{sec:setting}{1}, ``The recipe''):
\begin{enumerate}
  \item \emph{Build $K$.} A diagonal entry is how uncertain cohort $j$'s own
    estimate is around the portfolio level: its noise, averaged over its
    periods, plus $\tau^2$. An off-diagonal entry is how much two cohort
    estimates move together, which here happens only through shocks in the
    same period, so it is their averaged noise times $\rho_{ij}$
    (\hb{res:posdef}{7}).
  \item \emph{Check that $K$ is positive definite.}
  \item \emph{Read off the rest.} Invert $K$ once. The best portfolio
    weights, the variance of the portfolio estimate and every $z_j$ follow
    (\hb{res:gls}{3}, \hb{res:z}{5}).
\end{enumerate}

The point of the recipe is that steps 2 and 3 never change. Unequal series
lengths, a level shared within groups, or whatever comes next only change
how $K$ is built in step 1. So the theory is organized the same way, and so
is the Lean code: a \emph{core} proves everything that needs only a positive
definite $K$ and $\tau^2$, once; each \emph{model} builds its $K$ and proves
it positive definite; special cases are one-line corollaries. When a new
extension arrives, the first question is which $K$ it has.

\emph{Why positive definite, and not just semidefinite?} The thesis says
that $K$ is a covariance matrix and therefore positive definite
(\finding{1}). A covariance matrix is only positive \emph{semi}definite in
general, and the difference matters: step 3 inverts $K$. Take two cohorts
with no real difference between them ($\tau^2 = 0$) and identical,
perfectly correlated noise. Their estimates are then the same random number
twice; $K$ is singular, and there is no best way to weigh them because every
way is the same. In the model, what saves you is $\tau^2 > 0$: the noise part
of $K$ is semidefinite whenever the correlation table is, and adding $\tau^2$
to every diagonal entry lifts every direction strictly above zero. So the
two conditions to remember are $\tau^2 > 0$ and a positive semidefinite
$(\rho_{ij})$.

\emph{Why at least two cohorts?} With one cohort the portfolio \emph{is}
that cohort: its weight is one, the gap $\hmu_1 - \hmu$ is always zero, and
the formula for $z_1$ reads $0/0$ (\finding{2}). There is nothing to blend
with. The thesis never says so, but the reduction to Bühlmann--Straub needs
it.

\emph{With estimated parameters.} In practice you plug estimates into $K$,
and then positive definiteness is no longer automatic: the estimate of
$\tau^2$ is often negative, and a correlation table estimated entry by entry
need not be a correlation matrix. Truncate $\hat\tau^2$ at zero and shrink
the off-diagonal correlations a little towards zero, and the plug-in
$\hat K$ is positive definite again, whatever the data, as long as the
estimated table is semidefinite (\hb{res:damped}{8}).
Section~\ref{sec:estimation} goes further: a fit that needs no hand-set
shrinkage.

\watch
\begin{itemize}
  \item \emph{Check the correlation table before you invert.} An estimated
    $(\hat\rho_{ij})$ can have a negative eigenvalue. Shrinking towards zero
    does not repair that for every amount of shrinkage; repair the table
    first (for example, clip its negative eigenvalues).
  \item \emph{Positive definite is not the same as well conditioned.} Close
    to singular, the credibility factors swing hard with small changes in
    the inputs, and can leave $[0, 1]$ (Section~\ref{sec:z}).
  \item \emph{$K$ is about the estimates, not the observations.} Its
    entries depend on how each cohort's periods were averaged. Change the
    averaging weights and you change $K$; the best ones are in
    Section~\ref{sec:weights}.
\end{itemize}

\see \hbsec{sec:setting}{1}, \hb{res:gls}{3}, \hb{res:z}{5},
\hb{res:posdef}{7}, \hb{res:damped}{8} ·
\lean{Credibility.Model.K_apply},
\lean{Credibility.Model.K_posDef},
\lean{Credibility.Unequal.K_posDef},
\lean{Credibility.Unequal.dampedK_posDef},
\lean{Credibility.Core.bestZ_eq} ·
\app{Fitting steps} · \finding{1}, \finding{2} · thesis §2.3--2.5.

% ===========================================================================
\section{The weights}\label{sec:weights}

\oneline A cohort's estimate averages its periods by how precise each one
is, and the portfolio estimate averages the cohorts by the weights that make
it as precise as possible given $K$: the generalized least squares (GLS)
weights, which finance knows as the minimum-variance portfolio.

\idea Two averages come before any blending: a cohort's periods into
$\hmu_j$, and the cohorts into $\hmu$. Any weights that sum to one give an
unbiased estimate; the question is which give the most precise one.

\emph{Within a cohort}, weigh each period by the inverse of its noise
variance (\hb{res:invVar}{2}). A year with twice the volume has half the
noise variance and counts twice. With $\sigma_{jt}^2 = \sigma^2/w_{jt}$ this
is the volume-weighted mean, exactly as in Bühlmann--Straub. Correlation
between cohorts does not change it: within one cohort the noise of
different periods is uncorrelated, and how the cohort moves with the others
only shows up later, in $K$.

\emph{Across cohorts}, the variance of a weighted average
$\sum_j c_j\hmu_j$ is a quadratic form in $K$, and the weights that make it
smallest are $K^{-1}\mathbf 1$, scaled to sum to one
(\hb{res:gls}{3}). Call them $b_j$, and the variance they reach $s^2$.
Without correlation $K$ is diagonal and $b_j$ is proportional to
$1/K_{jj} = 1/(\tau^2 + v_j)$, where $v_j$ is the noise variance of
$\hmu_j$. Note the $\tau^2$: even a cohort measured without any noise
only tells you about its own level $\mu_j$, which sits a typical $\tau$ away
from $\mu$. That is why a huge cohort does not take over the portfolio, and
why Bühlmann--Straub weighs cohorts by their credibility factors and not by
their volumes (Section~\ref{sec:special}).

With correlation, a cohort that moves with many others gets \emph{less}
weight than its precision alone would earn: much of what it says, the
others already say. Weights can even turn negative. A negative weight is
the portfolio estimate using one cohort to cancel noise it shares with
another, the way a hedge works.

One consequence of optimality carries the rest of the theory: every cohort
covaries with the portfolio estimate by the same amount, $s^2$
(\hb{res:upsilon}{4}). If some cohort covaried more, moving weight away
from it would make $\hmu$ more precise, and the weights would not have been
the best. Section~\ref{sec:z} uses this to make the credibility factor
simple.

\emph{The finance connection.} Read the cohorts as assets and $K$ as the
covariance of their returns: $b$ is then the global minimum-variance
portfolio, the most studied object in portfolio theory. Its main lesson
carries over: weights computed from an \emph{estimated} $K$ are noisy, and
shrinking them pays when the shrinkage is chosen for the loss you actually
care about (Section~\ref{sec:estimation}; \file{studies/literature.md}, §1).

\watch
\begin{itemize}
  \item \emph{GLS weights are not exposure weights.} The portfolio estimate
    $\hmu$ is not total claims over total exposure, and the blended rates
    need not reproduce total claims. When that matters, see
    Section~\ref{sec:balance}.
  \item \emph{Negative and very large weights are fragile.} They are
    optimal for the true $K$, but they amplify errors in an estimated one.
  \item \emph{The weights are best for $\hmu$, and that is enough.} They are
    chosen for the portfolio level $\mu$, not for each blend; the
    credibility factor then makes the best of the $\hmu$ it is given. The
    moment layer confirms that $\hmu$ is best in true mean squared error
    too (Section~\ref{sec:moments}).
\end{itemize}

\see \hb{res:invVar}{2}, \hb{res:gls}{3}, \hb{res:upsilon}{4} ·
\lean{Credibility.invVarWeight_optimal},
\lean{Credibility.invVarWeight_value},
\lean{Credibility.glsWeight_optimal},
\lean{Credibility.glsWeight_value},
\lean{Credibility.mulVec_glsWeight},
\lean{Credibility.Core.Estimates.portEst_optimal} ·
\app{Models} · \file{studies/literature.md} §1 · thesis §2.5,
eqs.~(23)--(25).

% ===========================================================================
\section{The credibility factor}\label{sec:z}

\oneline The best credibility factor is the part of the gap between cohort
and portfolio that is real, divided by everything the gap varies by; with
correlation it can exceed one, and in extreme cases turn negative, without
anything having gone wrong.

\idea Take the regression reading of Section~\ref{sec:credibility}: $z_j$ is
the share of the gap $\hmu_j - \hmu$ you expect to be real. With $K$ and the
GLS weights in hand, both parts of that share are short
(\hb{res:z}{5}).
\begin{itemize}
  \item \emph{The real part} is $(1 - b_j)\tau^2$: the cohort's own effect,
    minus the share $b_j$ of it that already sits inside $\hmu$. A cohort
    that dominates the portfolio has little real gap left to show.
  \item \emph{The whole gap varies by} $\tau^2 + v_j - s^2$: the cohort's
    own variance minus the portfolio's. That it is this simple is the
    ``every cohort covaries equally with the portfolio'' fact of
    Section~\ref{sec:weights}.
\end{itemize}
So the recipe is: compute $K$, invert it once, and every $z_j$ follows.
Without correlation it collapses to $\tau^2/(\tau^2 + v_j)$, the
Bühlmann--Straub factor (\hb{res:diag}{6}).

\emph{How correlation moves it.} Two cohorts show it best. Cohort~1's factor
is Bühlmann--Straub's with its noise variance $\sigma_1^2$ replaced by
$\sigma_1^2 - \sigma_1\sigma_2\rho$ (\hb{res:two}{13}). The noise that cohort~1
shares with cohort~2 also sits in $\hmu$, so it cancels in the gap and does
not count against cohort~1. Positive correlation raises $z_1$, negative
lowers it, and unlike in Bühlmann--Straub, $z_1$ now depends on the other
cohort's size too.

\emph{Above one.} The thesis noticed that $z_j > 1$ can happen and proposed
clipping it to $[0, 1]$, without saying when it happens. For two cohorts the
answer is short: $z_1 > 1$ exactly when $\sigma_1 < \sigma_2\rho$, that is,
when cohort~1 is the precise one and the noise is correlated enough
(\finding{3}, \hb{res:z1_range}{14}). The condition does not involve $\tau^2$.
What happens is noise cancellation. The gap $\hmu_1 - \hmu$ is then mostly
cohort~2's noise with its sign flipped, and the best blend corrects for it
by moving \emph{past} $\hmu_1$, away from $\hmu$. Clipping throws that
correction away.

\emph{The singularity, and below zero.} Push the correlation further and
$z_1$ blows up, then comes back negative. In the thesis's example, cohort~1
has sixteen times cohort~2's volume and $\tau^2$ equals cohort~1's noise
variance; then $z_1 = 1$ at $\rho = 0.25$, $2.5$ at $\rho = 0.4$, $\pm 25$
on either side of $\rho = 0.5$, and $-1$ at $\rho = 0.75$. What goes on is
in the weights. As $\rho$ rises, cohort~1's GLS weight climbs towards one,
so $\hmu$ becomes cohort~1's own estimate and the gap shrinks to nothing.
At the singularity it is exactly zero; past it $b_1 > 1$, $\hmu$ overshoots
$\hmu_1$, and the gap changes sign. The factor $z_1$ flips sign \emph{with
the gap}, but the blend itself does not jump: on both sides it is $\hmu$
plus the same smoothly varying multiple of $\hmu_1 - \hmu_2$
(\finding{17}). Only at the singularity itself is the blend stuck at
$\hmu_1$, because every $z$ gives the same thing there. So a negative $z_1$
is the best blend written in an awkward coordinate, not an error. Reaching
the singularity at all takes very unequal cohorts and strong correlation
(\hb{res:z1_range}{14}).

\emph{Better numbers to report.} The awkwardness comes from measuring the
gap against a portfolio that contains the cohort itself. Two other numbers
avoid it (\finding{17}).
\begin{itemize}
  \item \emph{The own-data weight} $c_j = b_j + z_j(1 - b_j)$: the share of
    the blended rate that comes from cohort $j$'s own estimate, counting the
    part that reaches it through $\hmu$. It never blows up, for any number
    of cohorts: its distance from $b_j$ is at most $\tau^2$ over the smallest
    eigenvalue of $K$. In the example above it reads $1.000$, $1.038$,
    $1.064$, $1.070$, $1.154$ where $z_1$ read $1$, $2.5$, $25$, $-25$, $-1$.
    A value above one says what $z > 1$ was trying to say: the rate moves
    past the cohort's own figure.
  \item \emph{Two cohorts: the gap to the other cohort.} The blend is
    $\hmu$ plus $\tau^2/D$ times $\hmu_1 - \hmu_2$, where $D$ is the
    variance of that difference. This factor is always positive and smooth.
    With more cohorts the same trick (compare with the rest, reweighted)
    only moves the trouble to $b_j = 1$, so $c_j$ is the one to use.
\end{itemize}
Neither replaces $z_j$: it is the classical quantity, the results are
stated in it, and it means something else. In the same example at
$\rho = 0$, $z_1 = 0.5$ but $c_1 = 0.947$, because cohort~1 dominates
$\hmu$. Report $c_j$ next to $z_j$. And if you want to restrict anything,
restrict the blend, not $z_j$: clipping $z_j$ to $[0, 1]$ is what loses
accuracy here. With three or more cohorts $z_j$ is singular only when every
other weight is zero, so the sign flip is mostly a two-cohort affair, but
$z_j$ can still be large near it.

\watch
\begin{itemize}
  \item \emph{In a fitted model, look twice at extreme factors.} With the
    true parameters a large or negative $z_j$ is right. With estimated
    ones it more often means that $\hat K$ is close to singular, typically
    through $\hat\rho$ close to one or $\hat\tau^2$ close to zero
    (Section~\ref{sec:estimation}).
  \item \emph{$v_j$ is the noise of the cohort's estimate,} after averaging
    its periods, not of one period's observation.
  \item \emph{The same $z_j$ is best for next year's premium,} not only
    for the historical level (\hb{res:forecast}{19}).
\end{itemize}

\see \hb{res:z}{5}, \hb{res:diag}{6}, \hb{res:two}{13},
\hb{res:z1_range}{14}, \hb{res:forecast}{19} ·
\lean{Credibility.Core.bestZ_eq},
\lean{Credibility.Model.bestZ_eq},
\lean{Credibility.TwoCohorts.z1_eq},
\lean{Credibility.TwoCohorts.bestZ_mul_one_sub_b1},
\lean{Credibility.Core.abs_bestZ_mul_one_sub_le_eigenvalues},
\lean{Credibility.TwoCohorts.z1_gt_one_iff},
\lean{Credibility.TwoCohorts.z1_neg_iff},
\lean{Credibility.TwoCohorts.singularity_reachable_iff},
\lean{Credibility.Core.Estimates.bestZ_forecast_optimal} ·
\app{One draw} (the $z > 1$ warning) · \finding{3}, \finding{17} ·
thesis §2.5, eq.~(26); §2.7, eqs.~(32)--(34).

% ===========================================================================
\section{Special cases}\label{sec:special}

\oneline When $K$ is simple enough to invert by hand (no correlation, two
cohorts, or one covariance shared by all pairs), the credibility factors
have closed forms and the portfolio weights are the factors themselves,
rescaled to sum to one.

\idea \emph{Bühlmann--Straub.} No correlation, and noise variance
$\sigma^2$ over volume. Then $z_j = w_{j*}/(w_{j*} + \kappa)$ with
$w_{j*}$ the cohort's total volume and $\kappa = \sigma^2/\tau^2$, and the
portfolio weights are $b_j = z_j/z_*$ with $z_*$ the sum of the factors
(\hb{res:bs}{9}). This is the classical model, recovered exactly; it needs
at least two cohorts (\finding{2}). It is also balanced, a property
Section~\ref{sec:balance} takes up.

\emph{Two cohorts.} Bühlmann--Straub with $\sigma_1^2 - \sigma_1\sigma_2\rho$
in place of $\sigma_1^2$, as in Section~\ref{sec:z}, and again
$b_1 = z_1/(z_1 + z_2)$ (\hb{res:two}{13}). That second identity is not in
the published thesis; it appeared on the defence slides (\finding{4}). It
means that any split of a portfolio in two, such as the two levels of a
binary rating factor, keeps the classical weight formula under
correlation.

\emph{Equal covariances.} Suppose every pair of cohorts has the same noise
covariance $c$: a shock that hits all cohorts alike, such as inflation.
Then subtract $c$ from each cohort's noise variance and use the classical
formulas; the weights are again $z_j/z_*$ (\hb{res:equalcov}{15},
\finding{5}). Bühlmann--Straub is $c = 0$ and two cohorts are
$c = \sigma_1\sigma_2\rho$, so this one result contains the other two. The
reason in one line: equal covariances add the same amount to every entry of
$K$, a term that shifts every cohort alike and leaves the direction of the
best weights unchanged.

\emph{When do weights follow from factors?} In all three cases the
portfolio weights are the credibility factors rescaled. That is not true in
general. With three or more cohorts and unequal covariances it fails; in one
random example the weights are $(0.536, 0.226, 0.238)$ against
$z/z_* = (0.564, 0.177, 0.259)$ (\finding{4}). For exactly three cohorts (and no zero weight),
equal covariances are also necessary, a converse shown on paper but not in
Lean; from four cohorts on, other
structures work too (\finding{5}).

\watch
\begin{itemize}
  \item \emph{Equal correlations are not equal covariances.} If cohorts
    differ in size, the same $\rho$ for every pair gives different
    covariances $\sigma_i\sigma_j\rho$, and the shortcut fails already with
    three cohorts.
  \item \emph{The formulas look classical, but the inputs are not.} In the
    two-cohort and equal-covariance forms, a cohort's ``noise'' is reduced
    by what it shares with the others, so its factor depends on them.
  \item \emph{Forcing one covariance on unequal cohorts backfires.} As a
    fitted model, a single covariance for all pairs gives exploding factors
    when volumes differ; the app offers it only with that warning
    (\app{Models}, ``common'').
\end{itemize}

\see \hb{res:diag}{6}, \hb{res:bs}{9}, \hb{res:two}{13},
\hb{res:equalcov}{15} ·
\lean{Credibility.bs_z},
\lean{Credibility.Core.bestZ_of_diag},
\lean{Credibility.TwoCohorts.b1_eq_z},
\lean{Credibility.Core.b_eq_of_equalCov},
\lean{Credibility.Core.bestZ_of_equalCov},
\lean{Credibility.Model.bestZ_of_equalCov} ·
\app{Models} · \finding{2}, \finding{4}, \finding{5} · thesis §2.6, §2.7.


\oneline A blend is balanced when the blended rates, charged on the past
volumes, would have brought in exactly the claims that came in; every blend
is balanced for \emph{some} volume measure, Bühlmann--Straub happens to be
balanced for the exposures, and under correlation forcing exposure balance
costs accuracy, so if you must, shift all rates by the same amount.

\idea Practitioners like balance because it is easy to explain: the tariff
neither gives away nor overcharges the portfolio's own history. In
Bühlmann--Straub it comes for free. For any data, total premium at the
credibility rates equals total claims (\hb{res:bs}{9}). It is an identity,
not an expectation, so it holds whether or not the model is right, and it
survives cohorts observed over different periods (\finding{10}).

\emph{Why it holds there, and what it is in general.} Write the blend as
$\hmu_j^b = z_j\hmu_j + (1 - z_j)\hmu$. The blended rates are balanced for
volumes $u_j$ exactly when $u_j(1 - z_j)$ is proportional to $b_j$
(\hb{res:balance}{10}). So every blend is balanced for one volume measure,
$u_j \propto b_j/(1 - z_j)$, and the only question is whether that measure is
one you care about. In Bühlmann--Straub $b_j = z_j/z_*$ and
$1 - z_j = \kappa z_j/w_{j*}$, so the measure is the exposure $w_{j*}$: that
is the whole reason for the classical balance property. Change anything and
the measure moves. With region-specific noise and no correlation, it is the
precision $w_{j*}/\sigma_j^2$, not the exposure. With correlation some
$b_j$ are negative (Section~\ref{sec:weights}), so some balancing volumes
are negative: no positive volume measure balances the best blend.

\emph{Off-balance is a property of the optimum, not of the fit.} With the
true parameters the best blend already has a negative weight in every
correlated world of the app's study, and it misses exposure balance by more
than 1\% in half or more of the draws when cohorts differ in noise level
(study finding S27 in \file{studies/damping-theta.md}). So an off-balance of
a few percent is expected under correlation, not a sign that something went
wrong.

\emph{The cost of forcing it.} Suppose balance is a business requirement.
Three ways come to mind, and all of them cost accuracy.
\begin{itemize}
  \item \emph{Bend the factors.} Keep the weights $b_j$ and choose the
    factors among the balanced ones. The balanced blends form a line,
    $1 - z_j \propto b_j/u_j$, and the best point on it has a closed form
    for any loss weights and any correlation (\hb{res:bestBalanced}{11},
    \finding{15}). But the line fixes the \emph{shape} of the factors:
    large cohorts are pushed to $z_j \approx 1$ and small ones far from
    their best $z_j$, even below zero. In the app's study worlds it costs up
    to 96\% more error than the best blend.
  \item \emph{Recompute the collective} the Bühlmann--Straub way, as the
    $(1 - z_j)$-weighted mean of the cohort estimates. It is the worst of
    the three, and unstable when the weights $w_{j*}(1 - z_j)$ sum to little.
  \item \emph{Shift every rate by the same amount}, so that total premium
    matches total claims. This is benchmarking in small-area estimation
    (Wang, Fuller and Qu, 2008; \hb{res:commonShift}{12}). It is the
    cheapest of the three: at most 41\% more error in the same worlds,
    against 96\% for the best balanced factors, and nothing at all without
    correlation and with a common noise level.
\end{itemize}
These costs are exact, not simulated: with the true parameters every one of
these corrections is a linear rule, so its mean squared error is a formula
in the second moments (\file{studies/exact-balance.js}). They do depend on
the worlds chosen.

\emph{Why balance costs anything.} The exposure-weighted mean of the cohort
estimates carries this period's common shock in full. The GLS weights hedge
that shock away (Section~\ref{sec:weights}); balance puts it back into every
rate. Without correlation there is no common shock to hedge, and the cost
vanishes.

\emph{Is the shift the best balanced correction?} Among corrections that add
a share $c_jD$ of the imbalance $D$ to each blend, the common shift is
best for the volume-weighted error when cohorts are independent, for any
distribution with the given moments (\hb{res:commonShift}{12}). That is
Wang, Fuller and Qu's argument: the blend's errors are then uncorrelated
with the imbalance, so spreading it unevenly only adds noise. Under
correlation their premise fails, because our blend is the best linear
predictor only for independent cohorts, and the claim as stated is false
(\finding{16}). In a three-cohort counterexample the best shares are
$(0.177, 0.121, 0.227)$ instead of $1/6$ each. But the difference is small:
in 2000 random worlds the common shift was never more than 2.2\% worse than
the best shares, which would need the fitted covariances of the errors with
$D$. So the shift is the one to use.

\watch
\begin{itemize}
  \item \emph{Decide what balance is for before forcing it.} If the goal is
    accurate rates, do not force it: the best blend is balanced for its own
    measure $b_j/(1 - z_j)$. If total premium must match total claims,
    shift.
  \item \emph{Don't restore balance by bending $z$.} It is the costliest
    way and produces factors outside $[0, 1]$ for the small cohorts.
  \item \emph{Balance is not a test of the model.} It is an identity in
    Bühlmann--Straub and an expected miss under correlation; neither says
    whether the parameters are right.
  \item \emph{The app has no balance switch} for this reason; it reports
    the off-balance as a note, not a warning (\app{Many draws}).
\end{itemize}

\see \hb{res:bs}{9}, \hb{res:balance}{10}, \hb{res:bestBalanced}{11},
\hb{res:commonShift}{12} ·
\lean{Credibility.bs_balance},
\lean{Credibility.Unequal.balance_total},
\lean{Credibility.Core.balanced_iff},
\lean{Credibility.Core.balanced_iff_of_diag},
\lean{Credibility.Core.bestBalanced_optimal},
\lean{Credibility.Core.bestBalanced_cost},
\lean{Credibility.Core.Estimates.bestBalanced_mse_optimal},
\lean{Credibility.Core.bestCorrection_optimal},
\lean{Credibility.Core.imbalanceCov_of_diag},
\lean{Credibility.Core.Estimates.commonShift_mse_optimal_of_diag} ·
\app{Many draws} · \finding{10}, \finding{15}, \finding{16} ·
\file{studies/damping-theta.md} (S27--S29), \file{studies/vc-balance.js},
\file{studies/exact-balance.js} · \file{studies/literature.md} §4 ·
thesis §2.9.

% ===========================================================================
\section{Estimating the parameters}\label{sec:estimation}

\oneline The thesis estimators of the noise covariances and of $\tau^2$ are
unbiased for any distribution, given at least two periods and two cohorts;
plugged into $K$ they need repair or a better-structured fit, and the
fitted model's error is larger than its formula says, by a price of
estimation that can itself be estimated from the data.

\idea \emph{The unbiased estimators.} The noise covariance of two cohorts
is estimated from how their deviations from their own means move together
over the periods, and divided by a known factor that removes the bias, as
$T - 1$ does for an ordinary variance (\hb{res:estimation}{16}). On the
diagonal this is the Bühlmann--Straub $\hat\sigma_j^2$. The level variance
$\tau^2$ is the spread of the cohort means around the portfolio mean, minus
the part of that spread the noise explains. Both are unbiased with the
actual expectation, for any random variables with the model's moments
(Section~\ref{sec:moments}).

\emph{What they need.} At least two periods ($T \ge 2$) and two cohorts
($J \ge 2$); otherwise a denominator is zero (\finding{7}). The thesis does
not say so, but Bühlmann--Straub needs the same, so this is a gap in the
statement, not in the method. Its classical $\hat\tau^2$ is the uncorrelated
special case, unbiased too, once a typo is fixed: the thesis averages the
cohorts' $\hat\sigma_j^2$ with $1/T$ where it should be $1/J$
(\hb{res:bs_tau}{17}). With cohorts observed over different periods, the
$T - 2$ in the denominator becomes a volume-weighted count per pair
(Section~\ref{sec:unequal}).

\emph{The plug-in $\hat K$, and damping.} Unbiased does not mean usable.
$\hat\tau^2$ is often negative, and the table of estimated covariances need
not be positive semidefinite, so $\hat K$ may not be invertible. The first
repair is the one of Section~\ref{sec:K}: truncate $\hat\tau^2$ at zero and
multiply the off-diagonal correlations by a damping factor
$\theta < 1$; then $\hat K$ is positive definite whenever the estimated
correlation table is semidefinite (\hb{res:damped}{8}). How much to damp?
The best fixed $\theta$ grows with the correlation and with the length of
the history, from 0 without correlation to nearly 1 at $\rho = 0.9$ or 20
years, and no damping at all ($\theta = 1$) is risky with short
histories: with five years and eight regions, $\hat K$ fails to be positive
definite in nearly a quarter of the draws (study findings S1, S2).

\emph{Shrink for the loss you care about.} The obvious way to choose
$\theta$ is to make $\hat K$ as close as possible to the true $K$. That is
the wrong target. The rates want \emph{less} damping than the best estimate
of $K$ would get, most of all under heavy tails (S12). The reason, a
hypothesis so far: $\hat K$ is computed from the same years as the cohort
means, so a hard winter in this data set shows up in $\hat K$ and the GLS
weights cancel that very shock. $\hat K$ is worth more as a measurement of
\emph{this} history's common noise than as an estimate of the population
$K$ (S13). The minimum-variance portfolio literature reached the same
lesson in finance (Ledoit and Wolf 2017; \file{studies/literature.md} §1).
Leave-one-year-out cross-validation does not help either, for the same
reason: a held-out year's shock cannot be cancelled (S14).

\emph{A fit with no tuning.} Damping is a hand-set number, and the project's
rule is that a fitted model needs none. The fit the app uses instead,
\emph{correlated Bühlmann--Straub with credibility-weighted variances},
gets there by giving the covariances fewer parameters
(\file{studies/damping-theta.md}, ``Ansatz''):
\begin{itemize}
  \item \emph{One correlation} $\rho$ for all pairs, not one covariance:
    with unequal volumes one common covariance exceeds the noise variance of
    the large cohorts and their factors explode (S4; Section~\ref{sec:special}).
    With one $\rho \in [0, 1)$ the fitted $\hat K$ is positive definite by
    the theorem of Section~\ref{sec:K}, so damping disappears (S15).
  \item \emph{Variances by credibility.} Each cohort's $\hat\sigma_j^2$ is
    blended with the pooled one by a credibility factor whose own parameters
    (the noise of $\hat\sigma_j^2$, the spread of the true $\sigma_j^2$)
    come from moment equations: Bühlmann--Straub one level up (S17). A
    common $\sigma^2$ fails when regions really differ (S16), and a separate
    $\hat\sigma_j^2$ per region is fragile with short histories (S10).
  \item \emph{Every parameter by a moment equation}, nothing chosen by hand.
\end{itemize}
Over 45 simulated worlds it is the most robust of the fits tried: best or
near best almost everywhere, and equal to Bühlmann--Straub when there is no
correlation (\file{studies/variance-credibility.js}). The model, $K \succ 0$
and the blend are proved; the estimators' unbiasedness under this structure
is not yet in Lean, and the quality of the whole is numerical, as it is for
the EBLUP in small-area estimation (\file{studies/literature.md} §2).

\emph{The error of a fitted model.} Plug estimates into the
known-parameter error formula and you understate the error, by 13--50\% in
the study (S18): the formula misses the \emph{price of estimation}. Small-area
estimation splits the fitted model's error into the known-parameter error,
a part from estimating the parameters, and a cross term (Kackar and
Harville 1984, Prasad and Rao 1990; \file{studies/literature.md} §4). Under
normal noise the cross term vanishes; under heavy tails it is negative,
because a large shock inflates both a cohort's deviation and the variance
estimates, so the fitted factor shrinks hardest exactly when the raw mean
is most off (S21). That is how a fitted model can beat the best blend with
fixed weights (S9).

Two data-only estimates of the fitted error were tried. A jackknife that
re-fits the parameters with one year deleted (Jiang, Lahiri and Wan 2002) is
right on average within about $\pm15\%$ with ten years (S19), cautious under
heavy tails because it has no cross term (S23). Stein's unbiased risk
estimate, with the fitted covariances plugged in, is comparable, better
under heavy tails and too small with short histories (S25). Neither tells
you the error of \emph{this} data set: both estimate the average over data
sets like it (S26).

\watch
\begin{itemize}
  \item \emph{Two periods and two cohorts at least}, and more of both than
    that for anything stable.
  \item \emph{Check $\hat K$ before inverting it.} A negative $\hat\tau^2$
    or an indefinite correlation table is a sample artefact; truncate and
    repair, or use a fit whose $\hat K$ is positive definite by
    construction.
  \item \emph{Don't fit more covariances than the data carry.} Without
    correlation, plain Bühlmann--Straub is slightly better than any
    correlated fit (S6): that is the cost of estimating correlations that
    are not there. Feasible GLS has the same weakness.
  \item \emph{Report the fitted error as a typical error}, not as an
    interval around this data set's rates (\app{One draw}, ``Typ.\ err.'').
  \item \emph{Study numbers are numerical.} Findings S1--S26 come from
    simulated worlds; each is labelled in
    \file{studies/damping-theta.md} by how far it can be trusted if the
    world changes.
\end{itemize}

\see \hb{res:damped}{8}, \hb{res:estimation}{16}, \hb{res:bs_tau}{17},
\hb{res:integral_quad}{20} ·
\lean{Credibility.Model.sigRhoHat_unbiased},
\lean{Credibility.Model.tauHat_unbiased},
\lean{Credibility.Model.KHat_unbiased},
\lean{Credibility.Model.sigRhoHat_self},
\lean{Credibility.Model.tauHatBS_unbiased},
\lean{Credibility.Unequal.dampedK_posDef},
\lean{Credibility.Unequal.K_posDef} ·
\app{Models}, \app{Fitting steps} · \finding{7} ·
\file{studies/damping-theta.md} (S1--S26),
\file{studies/literature.md} §2, §4 · thesis §2.8, eqs.~(37)--(41).

% ===========================================================================
\section{Why the moments are real}\label{sec:moments}

\oneline Everything up to here treats covariances and mean squared errors as
given numbers; the moment layer proves that, for \emph{any} random
variables with those first and second moments, they are the actual
covariances, the blend is optimal in actual mean squared error, and such
random variables exist exactly when the correlation table is positive
semidefinite.

\idea \emph{Given numbers versus actual covariances.} The thesis, like most
credibility texts, works with second moments as symbols: $K_{ij}$ is
``the covariance of $\hmu_i$ and $\hmu_j$'', and the best $z_j$ minimises a
quadratic in those symbols. That is an algebraic optimisation. It becomes a
statement about real data only once someone checks that the symbols are
what they claim to be: that for random variables with the model's moments,
$\cov(\hmu_i, \hmu_j)$ actually equals the $K_{ij}$ of the formula, and the
quadratic actually equals $E(\hmu_j^b - \mu_j)^2$. In Lean that check is a
separate layer, and it uses no distribution: any probability space, any
square-integrable random variables with the stated means and covariances
(\finding{11}).

\emph{What the layer proves.}
\begin{itemize}
  \item \emph{The covariances are real.} $K$ is the actual covariance of the
    cohort estimates, and the portfolio estimate $\hmu$ is unbiased with the
    smallest actual variance among unbiased linear combinations
    (\hb{res:mse}{18}).
  \item \emph{The blend is optimal in true mean squared error.} The best
    $z_j$ of Section~\ref{sec:z} minimises $E(z\hmu_j + (1-z)\hmu - \mu_j)^2$,
    the actual expected squared error (\hb{res:mse}{18}). The same $z_j$ is
    best for next year's outcome, whose error is the level's plus next
    year's noise variance (\hb{res:forecast}{19}).
  \item \emph{``Unbiased'' means the actual expectation.} The expectation of
    any quadratic form in the data is computed from the moments alone, so
    the estimators of Section~\ref{sec:estimation} are unbiased whatever the
    distribution (\hb{res:integral_quad}{20}).
\end{itemize}
The layer is proved once, at the level of $K$: any unbiased estimates whose
covariances with each other and with the true levels are known. The
correlated cohorts, unequal lengths and group levels only supply those
covariances, so each is a short corollary.

\emph{The assumptions are consistent exactly when $\rho \succeq 0$.}
Section~\ref{sec:K} needed the correlation table to be positive
semidefinite for $K$ to be positive definite, a hypothesis the thesis did
not state (\finding{1}). The moment layer shows it is not an extra
assumption. Random variables with the model's moments exist if and only if
the table is positive semidefinite (\hb{res:exists}{21}). So whenever the
model describes anything at all, $K$ is positive definite for free. And if
an \emph{estimated} correlation table fails the test, no data-generating
process has those correlations: it is a sample artefact, to be repaired,
not a feature of the portfolio. With cohorts observed over different
periods the condition is weaker and again exactly right: per period, the
correlations of the cohorts observed then must be semidefinite; two cohorts
that never share a period can have any correlation (\finding{12}).

\emph{What ``best'' means.} Everything in this section is about the best
blend with \emph{fixed} weights: $z_j$ and $b_j$ computed from the true
moments, then applied to the data. A fitted model's weights depend on the
data, so it is a non-linear estimator and is not bound by that floor; under
heavy-tailed noise it can do slightly better (Section~\ref{sec:estimation}).
The app therefore calls the floor ``the best fixed blend''.

\watch
\begin{itemize}
  \item \emph{Distribution-free is not assumption-free.} The layer needs
    square-integrable random variables, uncorrelated periods, and cohort
    effects uncorrelated with the noise. A trend or a persistent shock
    breaks the second; noise that persists into next year belongs to the
    level, not to next year's noise.
  \item \emph{Second moments decide the blend, not its uncertainty.} How
    precise the \emph{estimated} parameters are depends on fourth moments,
    which the model leaves free; that is why the best damping and the
    fitted error depend on the noise distribution
    (\file{studies/damping-theta.md}, S7, S21).
\end{itemize}

\see \hb{res:mse}{18}, \hb{res:forecast}{19}, \hb{res:integral_quad}{20},
\hb{res:exists}{21} ·
\lean{Credibility.Core.Estimates.portEst_optimal},
\lean{Credibility.Core.Estimates.mse_blend},
\lean{Credibility.Core.Estimates.bestZ_mse_optimal},
\lean{Credibility.Core.Estimates.bestZ_forecast_optimal},
\lean{Credibility.Model.Noise.mse_blend},
\lean{Credibility.Model.Noise.bestZ_mse_optimal},
\lean{Credibility.Model.Noise.integral_quad},
\lean{Credibility.Model.integral_quad_of_cov},
\lean{Credibility.Model.exists_noise_iff},
\lean{Credibility.Unequal.Noise.bestZ_mse_optimal},
\lean{Credibility.Unequal.exists_noise_iff} ·
\app{Many draws} (fit against the floor) · \finding{1}, \finding{11},
\finding{12} · \file{studies/damping-theta.md} (S9).

\section{Cohorts observed over different periods}\label{sec:unequal}
\notyet{unequal lengths; the per-pair count.}

\section{Groups and shared levels}\label{sec:groups}
\notyet{a shared level drops out; several groups; what the estimators see;
towards hierarchical credibility.}

\section{Choosing the cohorts}\label{sec:cohorts}
\notyet{when merging pays off; equal credibility is not grouping; optimal
cohorts are intervals, and when they are not.}

\section{Open directions}\label{sec:open}
\notyet{crossed levels; time trends; the optional distribution layer.}

\section{Related work}\label{sec:related}
\notyet{credibility with common effects, small-area estimation, covariance
shrinkage: what is used from where.}

\appendix

\section{Changes from the thesis}\label{app:changes}

The thesis (Andblom, Stockholm University, 2023:2) holds up: every result
in its Section~2 is proved in Lean, some with a hypothesis it left unstated.
This appendix lists what changed, by thesis section, and then maps each
numbered thesis equation to its handbook result and Lean name. Equation
numbers refer to the published thesis.

\subsection*{Corrections and sharpenings}

\begin{description}
  \item[Notation (§2.3).] The cohort effect, $\varepsilon_j$ in the thesis,
    is written $\eta_j$ here, to keep it apart from the noise
    $\varepsilon_{jt}$ (Appendix~\ref{app:notation}).
  \item[$K$ is positive definite (§2.4, after eq.~(20)).] The thesis
    concludes this from $K$ being a covariance matrix, which only gives
    \emph{semi}definite. It holds because $\tau^2 > 0$ and the correlation
    table is positive semidefinite (\finding{1};
    Section~\ref{sec:K}), and the second condition is not an extra
    assumption: it holds whenever the model is consistent at all
    (\finding{11}; Section~\ref{sec:moments}).
  \item[At least two cohorts (§2.5--2.6).] With one cohort eq.~(26) reads
    $0/0$; the reduction to Bühlmann--Straub needs $J \ge 2$
    (\finding{2}).
  \item[Credibility above one (§2.7).] The thesis observes that $z_1 > 1$
    can happen. Now: below the singularity, exactly when
    $\sigma_1 < \sigma_2\rho$, whatever $\tau^2$ (\finding{3};
    Section~\ref{sec:z}).
  \item[Clipping is not needed (§2.7, eq.~(34)).] The thesis treats the
    sign flip of $z_1$ as a breakdown and clips $z_1$ to $[0, 1]$. With the
    true parameters the blend is smooth through the singularity, and a
    negative $z_1$ is still the best blend; clipping loses accuracy. Report
    the own-data weight $c_j$ next to $z_j$, which is bounded for any
    number of cohorts (\finding{17}; Section~\ref{sec:z}).
  \item[Weights from factors (§2.7.1).] The identity
    $b_1 = z_1/(z_1 + z_2)$ was on the defence slides, not in the published
    thesis. It holds for two cohorts, fails in general from three on, and
    extends to any number of cohorts with equal covariances (\finding{4},
    \finding{5}; Section~\ref{sec:special}).
  \item[Sample sizes for estimation (§2.8).] The unbiased estimators of
    eq.~(39) need $T \ge 2$ periods and $J \ge 2$ cohorts, which the thesis
    does not state (\finding{7}).
  \item[A typo in eq.~(13).] The Bühlmann--Straub $\hat\sigma^2$ averages
    the $J$ cohort estimates $\hat\sigma_j^2$, so the factor is $1/J$, not
    $1/T$ (\finding{7}; \hb{res:bs_tau}{17}).
  \item[Balance (§2.9).] The thesis shows that Bühlmann--Straub balances
    and that the correlated blend need not. Now: which volumes balance any
    given blend (\hb{res:balance}{10}); that the best blend constrained to
    balance is a poor way to get balance; and that a common shift of all
    rates is the best balanced correction without correlation, and close to
    it with (\finding{15}, \finding{16}; Section~\ref{sec:balance}).
\end{description}

\subsection*{Beyond the thesis}

New results, each explained in its own section: the moments are the actual
covariances of any random variables with them, so the blend is optimal in
true mean squared error, also as a forecast of next year
(Section~\ref{sec:moments}; \finding{11}, \finding{12}); cohorts observed
over different periods (Section~\ref{sec:unequal}; \finding{8}--\finding{10});
levels shared within groups (Section~\ref{sec:groups}; \finding{13}); and
choosing the cohorts: when merging pays, and when optimal cohorts are
intervals (Section~\ref{sec:cohorts}; \finding{6}, \finding{14}).

\subsection*{Concordance}

Each numbered equation of the thesis, with the handbook result that states
it now and the Lean names that prove it. Equations not listed are
definitions or steps of a derivation.

\begin{description}
  \item[(1)] The blend $\hmu_j^{\,b}$. \hbsec{sec:setting}{1}.
  \item[(2)--(5)] The three optimization problems: the weights within a
    cohort, the weights across cohorts, the factor. They can be solved one
    at a time (Section~\ref{sec:credibility}).
  \item[(6)--(7)] Error variance of the blend and the best factor.
    \hb{res:bestZ}{1} ·
    \lean{Credibility.BlendMoments.errVar_eq},
    \lean{Credibility.BlendMoments.bestZ_optimal},
    \lean{Credibility.BlendMoments.bestZ_unique}; in true mean squared
    error \lean{Credibility.Core.Estimates.mse_blend}.
  \item[(8)--(13)] The Bühlmann--Straub model and its estimators.
    \hb{res:bs}{9}, \hb{res:bs_tau}{17} ·
    \lean{Credibility.bs_z},
    \lean{Credibility.Model.tauHatBS_unbiased} (eq.~(13) with $1/J$).
  \item[(14)] The Parodi--Bonche factor: literature, not formalized.
  \item[(15)--(16)] The data model and the covariance of observations.
    \hbsec{sec:setting}{1}, \hb{res:exists}{21} ·
    \lean{Credibility.Model.exists_noise_iff},
    \lean{Credibility.Model.Noise.cov_obs}.
  \item[(17), (22)] Any linear estimates, and the factor for them.
    \lean{Credibility.Model.bestZ_linear}.
  \item[(18)--(21)] Noise covariances $\sigma_{ij}^2$, the portfolio
    variance and $K$. \hb{res:posdef}{7} ·
    \lean{Credibility.Unequal.K_eq_of_full},
    \lean{Credibility.Model.Noise.cov_cohortEst}.
  \item[(23)] Inverse-variance weights within a cohort.
    \hb{res:invVar}{2} ·
    \lean{Credibility.invVarWeight_optimal},
    \lean{Credibility.invVarWeight_value},
    \lean{Credibility.Model.Noise.cohortEst_optimal}.
  \item[(24)--(25)] GLS weights across cohorts and their variance $s^2$.
    \hb{res:gls}{3}, \hb{res:upsilon}{4} ·
    \lean{Credibility.glsWeight_optimal},
    \lean{Credibility.glsWeight_value},
    \lean{Credibility.mulVec_glsWeight},
    \lean{Credibility.Core.Estimates.portEst_optimal}.
  \item[(26)] The proposed credibility factor. \hb{res:z}{5} ·
    \lean{Credibility.Core.bestZ_eq},
    \lean{Credibility.Model.bestZ_eq},
    \lean{Credibility.Model.gap_eq},
    \lean{Credibility.Core.Estimates.bestZ_mse_optimal}.
  \item[(27)] The recipe in one display; $K_{ij} = \sigma_{ij}^2\rho_{ij} +
    \tau^2\delta_{ij}$ and it is positive definite. \hb{res:posdef}{7} ·
    \lean{Credibility.Model.K_apply},
    \lean{Credibility.Model.K_posDef}.
  \item[(28)--(30)] The Bühlmann--Straub special case. \hb{res:diag}{6},
    \hb{res:bs}{9} ·
    \lean{Credibility.bs_a},
    \lean{Credibility.bs_z},
    \lean{Credibility.Core.bestZ_of_diag}.
  \item[(31)--(32)] Two cohorts: the weights and the factor.
    \hb{res:two}{13} ·
    \lean{Credibility.TwoCohorts.gls_two},
    \lean{Credibility.TwoCohorts.z1_eq},
    \lean{Credibility.TwoCohorts.b1_eq_z}.
  \item[(33)--(36)] The singularity, clipping and the numerical example.
    \hb{res:z1_range}{14} ·
    \lean{Credibility.TwoCohorts.z1_gt_one_iff},
    \lean{Credibility.TwoCohorts.z1_neg_iff},
    \lean{Credibility.TwoCohorts.singularity_reachable_iff},
    \lean{Credibility.TwoCohorts.bestZ_mul_one_sub_b1}. The error curves
    (33), (35), (36) are plotted, not formalized.
  \item[(37)--(39)] The unbiased estimators of $\sigma_{ij}^2\rho_{ij}$,
    $\tau^2$ and $K$. \hb{res:estimation}{16} ·
    \lean{Credibility.Model.unbiased_dev},
    \lean{Credibility.Model.sigRhoHat_unbiased},
    \lean{Credibility.Model.tauHat_unbiased},
    \lean{Credibility.Model.KHat_unbiased}.
  \item[(40)] Truncating $\hat\tau^2$ at zero. \hb{res:damped}{8} ·
    \lean{Credibility.Unequal.dampedK_posDef}.
  \item[(41)] On the diagonal, the Bühlmann--Straub $\hat\sigma_j^2$.
    \hb{res:estimation}{16} ·
    \lean{Credibility.Model.sigRhoHat_self}.
  \item[§2.9] Balance (no equation number). \hb{res:bs}{9},
    \hb{res:balance}{10} ·
    \lean{Credibility.bs_balance},
    \lean{Credibility.Core.balanced_iff}.
\end{description}

\section{Concept index}\label{app:index}
\notyet{term, one-line meaning, section.}

\section{Notation}\label{app:notation}

The symbols of the paper, in the order they appear, with the
Bühlmann--Straub equivalent where there is one (as in Bühlmann and Gisler,
\emph{A Course in Credibility Theory}) and the thesis symbol where it
differs. The handbook uses the same notation (\hbsec{sec:setting}{1}).

\begin{center}
\small
\begin{tabular}{p{2.1cm}p{6.0cm}p{3.6cm}p{1.6cm}}
  \toprule
  Here & Meaning & Bühlmann--Straub & Thesis \\
  \midrule
  $j = 1, \dots, J$ & cohort (segment, rating class, card) & $i$, contract or risk & \\
  $t = 1, \dots, T$ & period (underwriting year) & $j$, year & \\
  $\mu_{jt}$ & observed key ratio of cohort $j$ in period $t$ & $X_{ij}$ & \\
  $w_{jt}$ & volume (exposure) of a cell & $w_{ij}$ & \\
  $w_{j*}$, $w_{**}$ & total volume of cohort $j$, of the portfolio & $w_{i\bullet}$, $w_{\bullet\bullet}$ & \\
  $\mu$ & overall level & $m$, collective premium & \\
  $\mu_j$ & true level of cohort $j$: the target & $\mu(\Theta_i)$ & \\
  $\eta_j$ & cohort effect, centred, unit variance & $\Theta_i$, a risk parameter & $\varepsilon_j$ \\
  $\tau^2$ & variance between cohorts & $a$ (Bühlmann--Gisler: $\tau^2$) & \\
  $\varepsilon_{jt}$ & noise, centred, unit variance & & \\
  $\sigma_{jt}$ & noise standard deviation of a cell & $\sigma/\sqrt{w_{ij}}$ & \\
  $\rho_{ij}$ & correlation of two cohorts' noise in the same period & $0$ for $i \ne j$ & \\
  $\delta_{ij}$ & 1 if $i = j$, else 0 & & \\
  $\kappa$ & $\sigma^2/\tau^2$ & $\kappa$ & \\
  \midrule
  $a_{jt}$ & weight of period $t$ in $\hmu_j$ & $w_{ij}/w_{i\bullet}$ & \\
  $\hmu_j$ & cohort estimate & $X_i$ & \\
  $\sigma_{ij}^2$ & covariance factor of the noise in $\hmu_i$, $\hmu_j$ & & \\
  $v_j$ & noise variance of $\hmu_j$ ($= \sigma_{jj}^2$) & $\sigma^2/w_{i\bullet}$ & $\sigma_{jj}^2$ \\
  $K$ & covariance matrix of the cohort estimates & diagonal & \\
  $\mathbf 1$ & vector of ones & & $\mathbf u$ \\
  $b_j$ & portfolio (GLS) weight & $\alpha_i/\alpha_\bullet$ & \\
  $\hmu$ & portfolio estimate & $\hat\mu_0$ & \\
  $s^2$ & variance of $\hmu$ & & \\
  $z_j$ & credibility factor & $\alpha_i$ & \\
  $z_*$ & sum of the factors & $\alpha_\bullet$ & \\
  $\hmu_j^{\,b}$ & the blend, the credibility estimate & $\widehat{\mu(\Theta_i)}$ & \\
  $c_j$ & own-data weight $b_j + z_j(1 - b_j)$ & $\alpha_i + (1 - \alpha_i)\alpha_i/\alpha_\bullet$ & \\
  \midrule
  $\sigma_1, \sigma_2$ & two cohorts: noise standard deviation of $\hmu_j$ & & \\
  $\kappa_j$ & two cohorts: $\sigma_j^2/\tau^2$ & & \\
  $\rho^*$ & two cohorts: where $z_1$ is singular & & \\
  $D$ & two cohorts: $\Var(\hmu_1 - \hmu_2)$ & & \\
  $T_j$ & the periods in which cohort $j$ is observed & & \\
  $L$ & covariance matrix of the true levels $\mu_j$ & $a\,I$ & $\tau^2 I$ \\
  $\varphi$ & share of $\tau^2$ shared within a group & & \\
  $u_j$ & volumes in a balance condition & $w_{i\bullet}$ & \\
  \bottomrule
\end{tabular}
\end{center}

\paragraph{Reading Bühlmann--Straub in this notation.} Bühlmann--Straub is
the case $\sigma_{jt}^2 = \sigma^2/w_{jt}$ and $\rho_{ij} = 0$ for
$i \ne j$. Of its independence assumptions exactly one is dropped: that two
cohorts' observations in the same period are uncorrelated
(Section~\ref{sec:model}). With it, $v_j = \sigma^2/w_{j*}$, $K$ is
diagonal, $z_j = w_{j*}/(w_{j*} + \kappa)$ and $b_j = z_j/z_*$
(\hb{res:bs}{9}). Two things to watch:
\begin{itemize}
  \item \emph{Two different $\kappa$.} The thesis's two-cohort $\kappa_j =
    \sigma_j^2/\tau^2$ is the noise variance of the cohort \emph{estimate}
    over $\tau^2$, so it already contains the volume and the number of
    periods; Bühlmann--Straub's $\kappa = \sigma^2/\tau^2$ is per unit of
    volume.
  \item \emph{$\sigma_{jt}$ is a standard deviation,} not a variance:
    $\sigma_{jt}^2 = \sigma^2/w_{jt}$.
\end{itemize}

\end{document}
